
\subsection*{Setting for $U_q(\mathfrak{sl}_2)$}

Let $k$ be a commutative ring with identity, and choose $q \in k^\times$ such
that $q-q^{-1}$ is invertible in $k$.
Let $U = U_q(\mathfrak{sl}_2)$ denote the quantized universal enveloping
algebra of $\mathfrak{sl}_2$ over $k$ with deformation parameter $q$. 
By following \cite{Ohtsuki}, It is generated by $E$, $F$, and $K^{\pm 1}$, subject to the relations

\begin{center}
    \renewcommand{\arraystretch}{1.4}
    \begin{tabular}{@{}c@{\qquad}c@{}}
        $KK^{-1} = K^{-1}K = 1$ & $EF - FE = \dfrac{K - K^{-1}}{q - q^{-1}}$ \\
        $KE = q^2 EK$ & $KF = q^{-2} FK$
    \end{tabular}
\end{center}
and
\begin{center}
    \renewcommand{\arraystretch}{1.4}
    \begin{tabular}{@{}c@{\qquad}c@{\qquad}c@{}}
        $\Delta(K) = K \otimes K$ & $\epsilon(K) = 1$ & $S(K) = K^{-1}$ \\
        $\Delta(E) = E \otimes K + 1 \otimes E$ & $\epsilon(E) = 0$ & $S(E) = -EK^{-1}$ \\
        $\Delta(F) = F \otimes 1 + K^{-1} \otimes F$ & $\epsilon(F) = 0$ & $S(F) = -KF$.
    \end{tabular}
\end{center}

\subsection*{Double-bosonization interpretation of $U$}
We here record Majid's interpretation on Lusztig's construction of $U$ \cite{Majid99}.

Let $H = k[K, K^{-1}]$ be a Hopf algebra with $\Delta(K) = K \otimes K$, $\epsilon(K) = 1$, and $S(K) = K^{-1}$, and let $A = k[t,t^{-1}]$ with another copy of $H$.
We let $A$ be dual quasitriangular with a univesal $r$-form $\mathscr{R}: A \otimes A \rightarrow k$ by letting 
\[
\mathscr{R}(t^i, t^j) = q^{2ij}.
\]

\begin{lemma}\cite[Lemma 4.1.]{Majid99}
    We have a pairing of Hopf algebras
    \[
    \langle \ ,\ \rangle: H \otimes A \rightarrow k, \quad \langle K, t \rangle = q^2
    \]
    and relative to it, we have a weak quasitriangular structure
    \[
    \mathscr{R}, \overline{\mathscr{R}}: A \rightarrow H, \qquad \mathscr{R}(t) = K, \quad \overline{\mathscr{R}}(t) = K^{-1}.
    \]
\end{lemma}
Note that here we have $\overline{\mathscr{R}} = \mathscr{R}^{-1}$.

Recall that a pair $(H,A)$ of Hopf algebras with a pairing called weak quasitriangular if there exists a convolution-invertible algebra and anti-coalgebra maps $\mathscr{R}, \overline{\mathscr{R}}: A \rightarrow H$ such that
\[
\langle \overline{\mathscr{R}}(a), b \rangle = \langle \mathscr{R}^{-1}(b), a \rangle, \quad a, b \in A
\]
and
\[
\partial^R h = \mathscr{R} * (\partial^L h) * \mathscr{R}^{-1}, \quad \partial^R h = \overline{\mathscr{R}} * (\partial^L h) * \overline{\mathscr{R}}, 
\]
where $(\partial^L h)(a) = \langle h_1, a \rangle h_2$ and $(\partial^R h)(a) = h_1 \langle h_2, a \rangle$ are left and right coregular actions of $H$.

Note that the above right $H$-action is generalized in the following sense:
Let $V$ be a left $A$-comodule.
One cane define a right $H$-action on $V$ by
\[
    v \triangleleft h = \langle h, v_{-1} \rangle v_0.
\]
With this convention, one can recognise that the braiding induced by dual quasitriangular structure of $A$ is equal to the braiding induced by weak quasitriangular structure of $(H, A)$.
Recall that the braiding $\Psi$ in the category of left $A$-comodueles is given by
\[
    \Psi(v \otimes w) = \mathscr{R}(w_{-1} \otimes v_{-1}) w_0 \otimes v_0,
\]
otherwise, the braiding induced from weak quasitriangular structure:
\[
    \Psi(v \otimes w) = w_0 \otimes (v \triangleleft \mathscr{R}(w_{-1})) = \langle \mathscr{R}(w_{-1}), v_{-1} \rangle w_0 \otimes v_0 .
\]
Both braidings are equal if the universal $r$-form is derived from the weak quasitriangular structure as follows:
\[
    \mathscr{R}(a \otimes b) = \langle \mathscr{R}(a), b \rangle.
\]

Now we let $B = k[E]$ be the free algebra generated by $E$ with a left $A$-comodule structure defined by $E \mapsto t \otimes E$, and induced right $A$-module structure defined by $E \triangleleft K = q^2 E$.
Moreover, if we let $E$ be a primitive element with $S(E) = -E$, one can show that $E$ is a braided Hopf algebra in the category of left $A$-comodules.

Dually, cosider $D = k[F]$ be the free algebra generated by $F$ with a right $A$-comodule structure $F \mapsto F \otimes t$, hence left $H$-module.
There exists a Hopf pairing between $B$ and $D$ defined by
\[
    \langle F, E \rangle = (q - q^{-1})^{-1},
\]
where the normalization factor is chosen by following Lusztig.
We also let $\overline{C} = D^{\text{cop}}$, which means that the (braided) comultiplciation of $C$ is set
\[
\Delta_{\overline{C}} = \Psi^{-1}_{D, D} \circ \Delta_D.
\]
We here regard that $\overline{C}$ is in the braided category of left $A$-comodules with $\overline{\mathscr{R}} = \mathscr{R}^{-1}$.

\begin{theorem}[\cite{Majid99}]
    There exists a unique Hopf algebra structure on $U = U(\overline{C}, H, B)$ on $\overline{C} \otimes H \otimes B$ such that two bosonizations $H \ltimes B$ and $\overline{C} \rtimes H$ are sub-Hopf algebras by the canonical inclusions, and with the cross relations.
\end{theorem}

Recall the cross relations in $H \ltimes B$ and $\overline{C} \rtimes H$ are given as
\[
    bh = h_1(b \triangleleft h_2), \quad hc = (h_1 \triangleright c)h_2
\]
which induces
\[
    EK = q^2 KE, \quad KF = q^{-2} FK.
\]
The cross relation of $U$ gives 
\[
[E, F] = (\mathscr{R}(t) - \overline{\mathscr{R}}(t))\langle F, E \rangle = \frac{K - K^{-1}}{q - q^{-1}}.
\]
The coproducts are given as follows:
\[
    \Delta(E) = E \otimes K + 1 \otimes E, \quad \Delta(F) = F \otimes 1 + K^{-1} \otimes K.
\]

Note that the substitution $q \mapsto q^{-1}$, $E \mapsto -E$ gives the usual convention of $U_q$.

In the Majid's convention, $U$ is quasitriangular with
\[
\mathscr{R} = \overline{\text{exp}}_B \mathscr{R}_H
\]
where
\[
    \sum f^\alpha \otimes \underline{S}(e_\alpha) = \mathrm{exp}_{q^{-2}}(-(q - q^{-1}) F \otimes E), \quad \mathscr{R}_H = q^{H \otimes H / 2}.
\]
Note that we are implicitly use the notation $K = q^H$.
In the usual convention,
\[
\mathscr{R} = \mathrm{exp}_{q^2}(-(q - q^{-1}) F \otimes E) \cdot q^{-H \otimes H/2}.
\]
Hence one can recognize that Majid uses $R_{21}^{-1}$ for Ohtsuki's $R$ (with a proper braided exponential convention).


\subsection*{Realization of the category of YD modules}

See the note \emph{Realization of YD category} in Notability.
Here, we consider the full subcategory of finite dimensional right-right Yetter--Drinfel'd modules.
By following \cite{Majid94}, that cateogory, with the fiber functor $F$ to the category of vector spaces, has an object $A$ satisftying the representability assumption for comodules, which is realized from the coend construction.
Hence, as there is an element in $\mathrm{Nat}(F, F \otimes H)$ representing the $H$-coaction, there exists a corresponding linear map $\phi_\delta \in \operatorname{Lin}(A, H)$.

Similarly, one can find a linear map $\phi_\mu \in \operatorname{Lin}(A, H^\circ)$ for the $H$-action as follows.
First, since every morphism in YD category is $H$-linear, there is also an element $\mu \in \mathrm{Nat}(F \otimes H, F)$ representing the $H$-action.
As all the objects have finite dimension, we have a transpose map $\mathrm{Nat}(F \otimes H, F) \rightarrow \mathrm{Nat}(F, F \otimes H^*)$.
This map is build that, for $\alpha \in \mathrm{Nat}(F \otimes H, F)$, and $X \in YD$, take a basis of $F(X)$, $\{ e_i \}$, and present
\[
\alpha(e_j \otimes h) = \sum  f^i_j (h) e_i
\]
and let $\hat{\alpha}(e_j):= \sum e_i \otimes f^i_j$.
Hence, $\hat{\mu}$ lies in $\mathrm{Nat}(F, F \otimes H^\circ)$, by the definition of $H^\circ$ related to the matrix coefficients.

\begin{question}
    Let $A$ be a Hopf algebra, and let $\phi_\delta: A \rightarrow H$ and $\phi_\mu: A \rightarrow H^\circ$ be Hopf algebra morphisms.
    For any $A$-comodules $V$ and $W$, we define $\Psi_{V,W}: V \otimes W \rightarrow W \otimes V$ defined as follows:
    \[
    \Psi_{V,W}(v \otimes w) := (w_0 \otimes v_0) \langle \phi_\mu(w_1), \phi_\delta(v_1) \rangle.
    \]
    Does $\{ \Psi_{V,W} \}$ satisfy the QYBE?
\end{question}